\section{Introducción}

El presente informe de implementación documenta el proceso de adopción, configuración y puesta en marcha del sistema de gestión en la nube basado en la plataforma Odoo para la empresa Machu Picchu Exclusive Tours E.I.R.L.. El objetivo central es digitalizar y optimizar los flujos críticos de gestión comercial y operativa, mejorando la seguridad de la información y la trazabilidad de los servicios ofrecidos.

Machu Picchu Exclusive Tours es una agencia de viajes cusqueña fundada en 2015, especializada en el sector de turismo receptivo de lujo y aventura con más de una década de experiencia; la organización se distingue por ofrecer experiencias auténticas lideradas por guías locales. Actualmente, la empresa opera de manera 100\% virtual, captando clientes exclusivamente extranjeros (principalmente de los mercados norteamericano y europeo) a través de recomendaciones y contactos directos por WhatsApp.

A pesar de su sólida reputación, la empresa enfrenta desafíos operativos derivados de una limitada digitalización. La gestión actual se apoya de forma manual en herramientas como Excel y WhatsApp, lo que provoca una dispersión de información confidencial ---incluyendo documentos de identidad y pasaportes de los pasajeros--- y dificulta el seguimiento de oportunidades comerciales y la trazabilidad de las reservas.

En este contexto, la implementación de Odoo se plantea como una solución estratégica para centralizar la data operativa, profesionalizar la atención al cliente y ordenar el flujo logístico comercial. Mediante el uso de módulos estándar como CRM y Proyectos, el sistema permitirá obtener una visión 360° del turista, gestionar itinerarios complejos de 2 a 20 días mediante tableros visuales y asegurar el control interno de adelantos y pagos antes de la compra de boletos no reembolsables.

El presente documento describe de manera estructurada las fases del proyecto, orientadas a transformar la gestión manual en un modelo digital inteligente. Este enfoque busca asegurar la autonomía de los datos, establecer reglas de negocio claras y dotar a la gerencia de herramientas analíticas para una toma de decisiones basada en evidencia.

